\documentclass[11pt]{article}

\usepackage[margin=1in]{geometry}
\usepackage{amssymb}
\usepackage{booktabs}
\usepackage{enumitem}
\usepackage{xcolor}
\usepackage{listings}
\usepackage[most]{tcolorbox}
\usepackage{hyperref}

\lstset{
  basicstyle=\ttfamily\small,
  columns=fullflexible,
  keepspaces=true,
  breaklines=true,
  frame=single,
  backgroundcolor=\color{gray!8},
  rulecolor=\color{gray!40},
  showstringspaces=false
}

\newtcolorbox{warnbox}{colback=red!5, colframe=red!60!black, fonttitle=\bfseries, title=Safety}
\newtcolorbox{tipbox}{colback=blue!4, colframe=blue!50!black, fonttitle=\bfseries, title=Tip}

\title{CodePath Intermediate Cybersecurity \\ Unit 1, Project 1 --- Catch Me if You Can! \\ \large Step-by-Step Guide and Submission Skeleton}
\author{}
\date{}

\begin{document}
\maketitle

\begin{warnbox}
The \texttt{.pcap} files contain \textbf{real Windows malware}. Do every step below inside your
\textbf{Ubuntu VM} (or on a Mac/Linux machine), never on the Windows host. Do not copy extracted
objects or \texttt{.eml} files back to Windows. Take your screenshots \emph{inside} the VM and move
only the \texttt{.png} images out.
\end{warnbox}

\section*{Deliverables at a glance}
\begin{enumerate}[label=\arabic*.]
  \item The malicious actor's \textbf{IP address}.
  \item The \textbf{subject lines} of three different phishing emails.
  \item A \textbf{detailed explanation} of how you found the actor.
  \item \emph{Stretch:} three \texttt{.eml} files (screenshots of them opened in a mail client).
\end{enumerate}

%----------------------------------------------------------------
\section{Set up the environment (in the VM)}
\begin{enumerate}[label=\textbf{1.\arabic*}]
  \item Boot your Ubuntu VM. Install Wireshark and the command-line tools:
\begin{lstlisting}
sudo apt update
sudo apt install -y wireshark tshark unzip
\end{lstlisting}
  \item Download \texttt{pcap\_files.zip} \emph{from inside the VM's browser} so it never touches Windows.
  \item Unzip it into its own folder and look at what you have:
\begin{lstlisting}
mkdir -p ~/pj1 && cd ~/pj1
unzip ~/Downloads/pcap_files.zip -d pcaps
ls -lh pcaps/
capinfos pcaps/*.pcap      # capture dates, packet counts, durations
\end{lstlisting}
\end{enumerate}

%----------------------------------------------------------------
\section{Triage: which capture has the phishing?}
The brief says only \textbf{one} capture contains malicious email. Find where the SMTP traffic is first.

\begin{enumerate}[label=\textbf{2.\arabic*}]
  \item Count SMTP packets per file:
\begin{lstlisting}
for f in pcaps/*.pcap; do
  echo "$f: $(tshark -r "$f" -Y smtp 2>/dev/null | wc -l) smtp packets"
done
\end{lstlisting}
  \item In the GUI, \textbf{Statistics $\rightarrow$ Protocol Hierarchy} gives the same picture visually.
  \item Keep notes on every file, including the ones that look clean. ``I ruled out file X because\ldots''
        is part of a good explanation.
\end{enumerate}

%----------------------------------------------------------------
\section{Read the emails}
\subsection*{Useful display filters}
\begin{center}
\begin{tabular}{@{}ll@{}}
\toprule
\textbf{Filter} & \textbf{What it shows} \\
\midrule
\texttt{smtp} & all SMTP commands and responses \\
\texttt{smtp.req.command == "MAIL"} & envelope sender (\texttt{MAIL FROM}) \\
\texttt{smtp.req.command == "RCPT"} & envelope recipients (\texttt{RCPT TO}) \\
\texttt{smtp.req.command == "DATA"} & start of each message body \\
\texttt{smtp.data.fragment} & the body data fragments \\
\texttt{imf} & reassembled messages (Internet Message Format) \\
\texttt{imf.subject contains "..."} & search by subject text \\
\bottomrule
\end{tabular}
\end{center}

\begin{enumerate}[label=\textbf{3.\arabic*}]
  \item Start with \texttt{imf} (or a \texttt{data} filter, as the hints suggest). Click a packet and expand the
        \emph{Internet Message Format} section in the details pane to see the headers.
  \item Right-click a packet $\rightarrow$ \textbf{Follow $\rightarrow$ TCP Stream} to read the entire SMTP
        conversation (\texttt{EHLO}, \texttt{MAIL FROM}, \texttt{RCPT TO}, \texttt{DATA}, headers, body).
  \item Dump a one-line-per-email summary to build your evidence table:
\begin{lstlisting}
tshark -r pcaps/<file>.pcap -Y imf -T fields -E header=y -E separator=' | ' \
  -e frame.number -e ip.src -e ip.dst -e imf.from -e imf.to -e imf.subject
\end{lstlisting}
  \item \textbf{Tip:} right-click a header field (e.g.\ \emph{Subject}) $\rightarrow$ \textbf{Apply as Column}
        to see it in the packet list.
\end{enumerate}

%----------------------------------------------------------------
\section{Decide: legitimate or fraudulent?}
For every email, record the fields below. BEC shows up in \emph{metadata}, not in obvious typos.

\begin{center}
\small
\begin{tabular}{@{}lllllll@{}}
\toprule
\textbf{Frame} & \textbf{Src IP} & \textbf{MAIL FROM} & \textbf{From:} & \textbf{Reply-To:} & \textbf{Subject} & \textbf{Verdict} \\
\midrule
\rule{0pt}{14pt} & & & & & & \\
\rule{0pt}{14pt} & & & & & & \\
\rule{0pt}{14pt} & & & & & & \\
\rule{0pt}{14pt} & & & & & & \\
\rule{0pt}{14pt} & & & & & & \\
\bottomrule
\end{tabular}
\end{center}

\noindent Questions to ask about each message:
\begin{itemize}
  \item Does the envelope sender (\texttt{MAIL FROM}) match the header \texttt{From:}?
  \item Is there a \texttt{Reply-To:} that points somewhere different from \texttt{From:}?
  \item Is the domain a look-alike (swapped letters, extra words, a different TLD)?
  \item Does the \textbf{source IP} match the IP that the organization's real mail comes from?
  \item What do the \texttt{Received:} headers say about the route the message took?
  \item Does the body ask for money, wire transfers, changed bank details, gift cards, urgency, or secrecy?
  \item Are there attachments (\texttt{Content-Disposition: attachment})? \emph{Do not open them.}
\end{itemize}

%----------------------------------------------------------------
\section{Attribute the actor}
\begin{enumerate}[label=\textbf{5.\arabic*}]
  \item Group your fraudulent emails by source IP. Do they share one sender?
  \item Confirm the legitimate emails come from a \emph{different} IP. That contrast is your strongest evidence.
  \item Check what else that IP did in the capture: \texttt{ip.addr == <suspect IP>}.
        \textbf{Statistics $\rightarrow$ Conversations} (IPv4 tab) helps here.
  \item Write down frame numbers for every claim, so a grader can verify each one.
\end{enumerate}

%----------------------------------------------------------------
\section{Stretch: export \texttt{.eml} files (VM only)}
\begin{enumerate}[label=\textbf{6.\arabic*}]
  \item With the phishing packets showing, go to \textbf{File $\rightarrow$ Export Objects $\rightarrow$ IMF\ldots}
  \item Save three phishing messages as \texttt{.eml} into a folder \emph{inside the VM}.
  \item Open each in a mail client in the VM (e.g.\ \texttt{sudo apt install thunderbird}) and screenshot it.
  \item Move only the screenshots out. If antivirus flags anything, that is expected: these are real samples.
\end{enumerate}

%----------------------------------------------------------------
\newpage
\section*{Submission skeleton}
\emph{Copy this structure into the Google Doc template, keeping the template's layout. Fill every
\texttt{[\ldots]} in your own words.}

\subsection*{Malicious actor IP}
\texttt{[IP address]}

\subsection*{Phishing email subject lines}
\begin{enumerate}
  \item \texttt{[subject 1]} \quad (frame \texttt{[\#]}, capture \texttt{[file]})
  \item \texttt{[subject 2]} \quad (frame \texttt{[\#]}, capture \texttt{[file]})
  \item \texttt{[subject 3]} \quad (frame \texttt{[\#]}, capture \texttt{[file]})
\end{enumerate}

\subsection*{How I found the malicious actor}
\begin{enumerate}[label=\textbf{\arabic*.}]
  \item \textbf{Triage.} \emph{[Which captures you checked, how, and why you ruled each one in or out.]}
  \item \textbf{Filtering.} \emph{[The filters you used and what each one revealed. Include screenshots.]}
  \item \textbf{Red flags.} \emph{[For each phishing email: the specific header or body evidence (From vs.\ Reply-To, domain, request).]}
  \item \textbf{Attribution.} \emph{[How the source IP ties the emails together, and how it differs from legitimate mail.]}
  \item \textbf{Why metadata matters.} \emph{[One or two sentences connecting your evidence to the Key Takeaways.]}
\end{enumerate}

\subsection*{Stretch: \texttt{.eml} screenshots}
\texttt{[Screenshot 1]} \quad \texttt{[Screenshot 2]} \quad \texttt{[Screenshot 3]}

\subsection*{Submission checklist}
\begin{itemize}[label=$\square$]
  \item All required challenges complete
  \item Template layout followed; checklist in the doc ticked
  \item Images display correctly
  \item Doc sharing set to ``Anyone with the link can Edit''
\end{itemize}

\end{document}
